%*****************************************************************************************************            
%                               INSTITUT MINES-TELECOM, IMT ATlantique
%                         Département Mathematical & Electrical Engineering           
%                                     Technopôle de Brest-Iroise                   
%                                  CS 83818 - 29238 BREST Cedex 3                       
%            
%                                    **  Tous droits réservés **
%                                      
%
% NOM DU PROJET : TP LaTeX 006 / Modèle de rapport de recherche (validé par la DRI), de support de cours
% et de TP au format LaTeX avec les éléments de la charte graphique de la communication d'IMT Atlantique
%
% LIVRABLE : [ ] Memo, [ ] Résumé, [ ] Article, [ ] Code, [X] Rapport , [ ] Thèse, [ ] Cours/TP
%
% AUTEUR(s) : Thierry LE GALL (Dépt. MEE), Guillaume ANSEL (SC), Thomas GUILMENT (Dépt. SC)
%
% HISTORIQUE :
%
% Date          Nom(s)          Version  Description
% ------------  --------------- -------  ----------------------------------------------------------
% 12-mar-2015   T. LE GALL (SC) 0.0      création du document 
% 26-mar-2016   T. LE GALL (SC) 0.1      mise à jour : essai du package UTF8
% 26-sep-2016   G. ANSEL        0.2      mise à jour : ajout exemple utilisation TBannexe
%                                        ajout correction titres en majuscule dans sommaire
%                                        (cf BUGFIX_TITRES_SOMMAIRE dans imta.sty)
%                                        changement polices utilisées
% 09-jan 2017   T. LE GALL (SC) 0.3      mise à jour : évolution du modèle de TB vers IMT Atlantique
% 03-fev 2017   G. ANSEL        0.4      mise à jour : modifications pour adaptation à version 1.5 du package imta
%                                        ajout méta-données PDF
% 14-fev-2017   T. GUILMENT     0.5      mise à jour : ajout package fancyhdr pour entêtes et pieds de page
%                                        changement style titres
% 16-fev-2017   G. ANSEL        0.6      mise à jour : modifications pour version 1.6 package imta
%                                        ajout style de page IMTAfancy pour entêtes et pieds de page
%                                        remplacement \listoffigures et \listoftables par \IMTAlistefigures
%                                        et \IMTAlistetableaux
% 17-fev-2017   G. ANSEL        0.7      mise à jour : déplacement définition style IMTAfancy dans imta.sty
%                                        définition commande \IMTAfooter pour personnaliser le pied de page
%                                        en bas à gauche
% 21-fev-2017   T. LE GALL (SC) 0.8      mise à jour : ajout d'un exemple de n° de contrat de recherche
%                                        (pied de page) et vérification de la compilation du document 
% 28-fev-2017   G. ANSEL        0.9      mise à jour : ajout exemple d'utilisation de l'option copyright
%                                        du package imta (informations copyright sur quatrième de couverture)
% 06-fev-2017   T. LE GALL (SC) 1.0      mise à jour : ajout de définitions d'ensembles mathématiques et exemple
%                                        d'édition de bloc dans le texte.  Ajout de la mention FIN DE FICHIER. 
%                                        Test de compilation avec imta.sty (v2.3) = PASS (0 errors, 0 warnings).
% 06-avr-2017   T. LE GALL (SC) 1.1      mise à jour : ajout du type de document et du numéro de rapport
%                                        de recherche en tant que paramètre du document et affichage des
%                                        informations en première de couverture (suite à demande de la DRI)
% 03-aoû-2017   T. LE GALL (SC) 1.2      mise à jour : exemple de section de texte sur deux colonnes et équations
%                                        dans deux blocs adjacents de style Beamer
% 09-aoû-2017   T. LE GALL (SC) 1.3      mise à jour : ajout des équations (matrices, vecteurs)
% 11-sep-2017   T. LE GALL (SC) 1.4      mise à jour : attribution du n°ISSN 2556-5060 pour les rapports de recherche
% 22-sep-2017   G. ANSEL        1.5      mise à jour : ajout d'un exemple d'insertion de tableau
% 25-sep-2017   T. LE GALL      1.6      mise à jour : légende du tableau et n° de version du modèle en page de garde
%                                        + tests de non-négression du modèle (env. Win 7) suite à BUGFIX_21092017_1
%                                         et BUGFIX_21092017_2 dans le fichier imta.sty v2.8 et v2.9
% 27-nov-2018   G. ANSEL        1.7      mise à jour : modifications pour passage en version 3.5 du package imta
% 21-jan-2019   T. LE GALL (SC) 1.8      mise à jour : couverture : type de document au-dessus de titre et liste
%                                        des auteurs et contributeurs sous le titre + tests de compilation (Win10)
%                                        + tests de compilation (UBUNTU 18.04 LTS)
% 10-mai-2020  T. LE GALL (SC)  1.9      mise à jour : ajout d'exemples supplémentaires (équations)
%
% 25-jan-2021  T. LE GALL (MEE) 2.0      mise à jour : département Mathematical & Electrical Engineering
%
% 18-oct-2023  T. LE GALL (MEE) 2.2      mise à jour : 1ere + 4eme de courverture suite à demandes DRI (N. Fontaine)
%                                        - 1ere de couvertture: nouveau format et procédure de n°de rapport de recherche
%                                        - 1ere de couvertture: mentions des niveaux de classification suivant ZRR
%                                        - 4ème de couverture:
%                                          suppression site de Toulouse + maj.carte
%                                          suppression ISSN remplacé par CreativeComms
%                                          ajouts des icones cliquables (réseaux sociaux)                           
%
% 28-fév-2025  T. LE GALL (MEE),
%              J-L. CASTAIGNE (CARAE) 2.3 mise à jour : changement par la DIRCOM} des spécifications de la tonalité
%                                          de la couleur verte , codage passé de 164,210,51 à 153,204,153
%
% NOTES/COMMENTAIRES :
%
% - Exemple de rapport scientifique au format IMT Atlantique sous LaTex avec insertion de formules,
%   figures, ref. biblio, hyperliens, signets, renvois actifs, surlignage des formules. Le fichier
%   d'édition tp_latex_006.tex utilise le fichier de style imta.sty. Cf. fichier lab_006_readme.txt
%   situé sous le répertoire lab_006.
%
% REFERENCES :
%
% - cf. bibliographie
% 
%***********************************************************************************************

%-----------------------------------------------------------------------------------------------
%                                    DEBUT DU FICHIER D'EDITION
%-----------------------------------------------------------------------------------------------




\documentclass{article}
\usepackage{graphicx}
\usepackage{pdfpages}
\usepackage{hyperref}
\usepackage{subcaption}
\usepackage{float}
\usepackage{tabularx}
\usepackage{pdflscape}
\usepackage[table]{xcolor}
\usepackage{titling}

% Resserrement propre de la page de titre
\setlength{\droptitle}{-2.5em}

\title{Synthèse du projet v1}
\author{%
  Gautier de Marsac \\
  Maëlle Le Devedec \\ 
  Marc Duboc \\ 
  Maël Petit \\ 
  Raphaël Lebel \\ 
  Loann Nicotte
}
\date{%
  Smart Chess — Projet n°38\\[0.3em]
  \small Relecteur : Matthieu Arzel \quad • \quad Client : Pierre Tremenbert \\
   02/10/2025
}



\begin{document}

\maketitle
\centering
\includegraphics[width=5cm]{IMTA logo.png}
\newpage
\tableofcontents
\newpage

\raggedright

\section{Introduction}

Dans le cadre de l'UE Projet Commande Entreprise, nous avons décidé de proposer notre propre sujet aux responsables de l'UE: il s'intitule Smart Chess. Ce projet part du constat suivant: pour un joueur qui n'a pas l'habitude de jouer avec des échiquier virtuels sur ordinateur ou smartphone, il peut être difficile de trouver un partenaire de son niveau pour jouer et donc progresser. Pour remédier à cela, nous avons eu l'idée de concevoir un échiquier intéractif doté d'une intelligence artificielle capable de simuler un adversaire ou de suggérer des coups à un joueur en difficulté dans une partie: c'est le Smart Chess. 

\section{Objectifs}

Concrètement, ce projet est divisé en deux équipes qui avancent parallèlement avant de se rejoindre à l'issue du projet. La première équipe, constituée de Maëlle, Maël, Raphaël et Gautier, est en charge de l'intelligence artificielle. Ses objectifs sont de concevoir un chess bot capable d'une part de jouer une partie contre un adversaire et d'autre part de suggérer des coups à un joueur en difficulté selon des paliers (assistance plus ou moins élevée en fonction du palier sélectionné). \\
La deuxième équipe, composée de Marc et Loann, est en charge de la conception du plateau d'échecs. L'objectif est que l'échiquier reconnaisse le mouvement de pièces et illumine les cases pour les suggestions de coups. A l'issue du projet, nous envisageons également d'intégrer un système pour que les pièces controlées par l'IA se déplacent de manière autonome. \\
Avec ce projet, nous prévoyons que les personnes isolées pourront s'entraîner et progresser seules à tout moment de la journée avec des sensations et des repères similaires à une partie à deux joueurs sur un échiquier physique. 

\section{Diagramme de Gantt}

\begin{landscape} 
\begin{figure}[!htbp]
    \centering
    \includegraphics[width=\linewidth]{GanttPhase.png}
    \caption{Diagramme de phase}
\end{figure}

\vspace{1em} % espace entre les images

\begin{figure}[!htbp]
    \centering
    \includegraphics[width=\linewidth]{GanttJalon.png}
    \caption{Diagramme des Jalons}
\end{figure}
\end{landscape}

\section{Tableau des livrables}
\begin{figure}[!htbp]
    \centering
    \includegraphics[width=\linewidth]{TabLiv.png}
    \caption{Diagramme de phase}
\end{figure}

\section{Organigramme des tâches (WBS)}

\begin{figure}[H]
    \centering
    \includegraphics[width=0.8\linewidth]{Capture d’écran 2025-10-02 à 10.10.24.png}
    \caption{Organigramme WBS}
    \label{fig:placeholder}
\end{figure}

\section{Moyens de communication}

L’équipe communique en interne via un groupe WhatsApp pour les échanges rapides et de deux groupes WhatsApp dédiés a la communication aux seins des sous-équipes. Les deux sous équipes se réunissent également deux fois par semaine pour assurer le suivi du projet. La communication avec l'encadrant se fait sur Discord. Une réunion toute les deux semaines est programmé pour échanger avec le client, de même avec l'encadrant. Le partage de documents se fait sur Google Drive, les comptes rendus de réunions sont présents et accessible au client sur ce drive pour qu'il puisse suivre l'avancement de façon autonome. De plus les échanges par mail avec le client sont géré par Marc, responsable désigné. Enfin le code est versionné sur GitHub.

\section{Moyens de reporting}
CR un jeudi sur 2 au client et à l'encadrant technique pour faire le point et présenter nos problèmes.
\section{Méthodologie }
Approche itérative par prototypage et cycles de 2 semaines. Emploi des méthodes agiles pour une efficacité accrue.
\section{Positionnement stratégique}
\subsection{Public visé}
\subsubsection{Profil démographique des joueurs d’échecs}

\begin{itemize}
    \item Environ \textbf{70 000 joueurs licenciés en 2025} (contre 54 000 en 2018), avec une croissance régulière des effectifs, notamment chez les jeunes et les femmes (qui représentent environ 20 \% des joueurs).
    \item Le public regroupe aussi bien des \textbf{juniors} (moins de 18 ans) que des \textbf{seniors}, avec une forte présence des \textbf{18-35 ans}, souvent attirés par les aspects compétitifs et technologiques.
    \item Le marché des échecs (compétitions, plateaux, abonnements\dots) est également en augmentation : évalué à \textbf{3,4 milliards USD en 2024}, il passe à \textbf{3,70 milliards USD en 2025}, et les prévisions l’estiment à \textbf{7,64 milliards USD en 2032}.
\end{itemize}

\textit{Source : \url{https://www.fortunebusinessinsights.com/fr/chess-market-113098}}

\subsubsection{État du marché des robots de divertissement}

\begin{itemize}
    \item Le marché des robots de divertissement est évalué à \textbf{3,1 milliards USD en 2023}, à \textbf{3,86 milliards USD en 2024}, et devrait atteindre \textbf{18,1 milliards USD en 2032}, avec un taux de croissance annuel composé de \textbf{21,3 \%} sur la période 2024--2032.
    
    \item La croissance est fortement tirée par la demande croissante de \textbf{robots éducatifs} auprès des enseignants, des établissements scolaires et des parents. Ces robots sont utilisés pour développer des compétences en résolution de problèmes, en pensée critique et en créativité, en plus de leur dimension ludique. Des initiatives sont en cours, telles que le programme \textit{Build Your First Robot} en Inde.
    
    \item Les principaux utilisateurs finaux se répartissent entre les \textbf{médias, l’éducation, la vente au détail et autres secteurs}. En 2022, le segment de la \textbf{vente au détail} a dominé le marché, notamment grâce à la forte demande de jouets robotiques pour enfants.
\end{itemize}

\begin{figure}[h]
    \centering
    \begin{subfigure}[b]{0.3\linewidth}
        \centering
        \includegraphics[width=\linewidth]{MarchNum.png}
        \caption{Marché total (milliards USD)}
    \end{subfigure}
    \hfill
    \begin{subfigure}[b]{0.3\linewidth}
        \centering
        \includegraphics[width=\linewidth]{TypeMarche.png}
        \caption{Marché  par utilisateur (milliards USD)}
    \end{subfigure}
    \hfill
    \begin{subfigure}[b]{0.3\linewidth}
        \centering
        \includegraphics[width=\linewidth]{RegMarche}
        \caption{Part du marché dans le monde (milliards USD)}
    \end{subfigure}
\end{figure}
\textit{Source : \url{https://www.marketresearchfuture.com/fr/reports/entertainment-robots-market-2925 }}

\subsubsection{Analyse de ces résultats}

Le marché des échecs est toujours attractif au vu des résultats énoncés plus haut. Le public visé se compose surtout de personnes entre \textbf{18 et 35 ans}, mais aussi de \textbf{parents} qui recherchent des solutions éducatives pour développer les compétences de leurs enfants. De plus, au vu de la présence de niveaux de difficulté, d’indices pour gagner et d’analyses d’après-match, toute personne, qu’elle soit \textbf{novice ou compétitrice}, pourrait être intéressée par notre produit.

Les \textbf{clubs d’échecs} et les \textbf{établissements scolaires} représentent une cible stratégique, car ils pourraient intégrer ce plateau comme outil pédagogique moderne. Comme on l’a vu, plusieurs initiatives du gouvernement ont été faites dans ce sens. Une des stratégies de commercialisation serait de développer des \textbf{partenariats institutionnels}, mais cela nécessiterait de garder un \textbf{prix abordable}.

Enfin, les \textbf{échecs comme les robots de divertissement} suscitent un intérêt dans de nombreux pays et sur plusieurs continents. Si le produit se vend bien en France, une \textbf{commercialisation internationale} est plus qu’envisageable.

\subsection{Reformulation du besoin}

L’objectif du projet est de concevoir un \textbf{échiquier électronique en bois} capable d’assister les joueurs durant leurs parties.  
Cet échiquier répond à trois finalités principales :

\begin{enumerate}
    \item \textbf{Équilibrer les parties entre joueurs de niveaux différents} \\
    Grâce à une assistance intelligente (conseils de coups, correction d’erreurs, adaptation du niveau de difficulté), l’échiquier permet à un joueur débutant de jouer à armes égales contre un joueur expérimenté.  
    Cette fonction favorise à la fois le plaisir de jeu et l’apprentissage progressif des échecs.
    
    \item \textbf{Apporter une dimension pédagogique et technique} \\
    \begin{itemize}
        \item Sur le plan \textit{échiquéen}, l’échiquier accompagne le joueur dans sa progression et l’aide à développer sa réflexion stratégique.
        \item Sur le plan \textit{technologique}, le système constitue un support éducatif :
        \begin{itemize}
            \item une \textbf{IA simple et open source} permet de s’initier au codage et à la logique algorithmique ;
            \item l’échiquier est livré en \textbf{kit}, offrant la possibilité d’un montage manuel pour développer des compétences pratiques en électronique, mécanique et programmation.
        \end{itemize}
    \end{itemize}
    
    \item \textbf{Allier modernité et esthétique} \\
    L’échiquier conserve l’aspect traditionnel et noble du bois, gage de durabilité et d’attrait visuel.  
    L’intégration de la technologie (déplacement automatique des pièces, assistance IA) est réalisée de manière discrète et harmonieuse afin de préserver l’expérience authentique du jeu.
\end{enumerate}

\bigskip
\noindent
\textbf{En résumé}, le Smart Chess est à la fois :
\begin{itemize}
    \item un \textbf{outil d’apprentissage} des échecs et des technologies associées,
    \item un \textbf{moyen d’équilibrer les parties} entre joueurs de niveaux différents,
    \item et un \textbf{objet esthétique et innovant}, combinant tradition et modernité.
\end{itemize}

\subsection{Analyse des risques}

Le tableau ci-dessous présente les principaux risques identifiés pour le projet \textit{Smart Chess}, ainsi que les mesures de prévention et de remédiation envisagées.

\begin{table}[H]
\centering
\renewcommand{\arraystretch}{1.3}
\begin{tabularx}{\textwidth}{|X|X|X|}
\hline
\textbf{Risque} & \textbf{Prévention} & \textbf{Remédiation} \\
\hline
Blocage ou imprécision du mécanisme de déplacement des pièces & Tests précoces du système mécanique, choix de moteurs fiables & Ajustements mécaniques, remplacement des pièces défectueuses \\
\hline
Mauvaise détection des mouvements des pièces par le joueur & Calibration fine des capteurs, tests en conditions réelles & Correction logicielle, recalibrage automatique \\
\hline
Intégration IA–plateau difficile (latence, bugs) & Prototypage rapide de la communication IA–hardware & Débogage et mise à jour logicielle \\
\hline
Dépassement du budget dû au coût des composants (bois, moteurs, capteurs) & Recherche de fournisseurs alternatifs, estimation budgétaire précise & Remplacer certains matériaux par des équivalents moins coûteux \\
\hline
Disponibilité limitée de composants électroniques (pénuries) & Anticiper les commandes, prévoir des références alternatives & Redéfinition partielle de la conception matérielle \\
\hline
Complexité du montage en kit pour les utilisateurs & Fournir un manuel clair, tutoriels vidéo & Assistance en ligne, simplification du design \\
\hline
IA trop faible ou trop complexe pour les débutants & Ajustement des niveaux ELO, tests utilisateurs & Mise à jour de l’IA avec de nouveaux paramètres \\
\hline
IA insuffisamment puissante face à des joueurs expérimentés & Évaluer plusieurs moteurs d’échecs existants, choisir une base robuste et adaptée & Intégrer un moteur plus avancé (ex. Stockfish allégé) ou permettre la connexion à un serveur externe \\
\hline
Concurrence forte sur le marché & Positionnement prix clair, différenciation par le kit et l’aspect éducatif & Repositionnement marketing, ajout de fonctionnalités \\
\hline
\end{tabularx}
\caption{Analyse des risques du projet Smart Chess (AMDEC simplifiée, avec risque lié à la puissance de l’IA).}
\end{table}



\subsection{Concurence}
Sur le marché des échecs connectés, la majorité des concurrents proposent des caractéristiques communes : la plupart intègrent une IA réglable ou adaptative, permettant de jouer à différents niveaux, et proposent des indicateurs lumineux pour guider le joueur sur les coups possibles ou signaler les erreurs. Beaucoup sont accompagnés d’une application ou d’un logiciel compagnon, offrant des modes de jeu variés : face à face, en ligne ou contre l’IA. 
Chacun présente néanmoins des spécificités propres :
\begin{itemize}
    \item Particula (Israël) – 292-410 € : met l’accent sur ses caractéristiques compact et portable pour plus de mobilité. L’application gratuite propose trois modes de jeu et la difficulté de l’IA est réglable.
    \item Chessnut (Chine) – 737-935 € : il mise sur son système vocal et son moteur d’échecs intégré, avec une IA réaliste type Maia et une connectivité extensible.
    \item ChessUp (États-Unis) – 350-490 € : son point fort est l’écran tactile couleur et le Wi-Fi intégré pour les mises à jour. Il permet de jouer en ligne, d’accéder à des leçons et de régler l’assistance de l’IA pour équilibrer les parties.
    \item Sense Robot (Chine)– 1170 € : se distingue par son bras robotisé, ses 25 niveaux d’IA, plus de 1200 exercices, un entraîneur de finales, le coaching vocal et l’intégration complète avec Lichess.
    \item Miko (Inde)– 550-303 € : propose une IA adaptative, la réinitialisation automatique des pièces, la diffusion de matchs professionnels en direct et un guidage LED pour les coups et les erreurs. Il permet de défier des adversaires du monde entier et propose un coaching personnalisé.
\end{itemize}
Il ne sera donc pas facile de se faire une place sur ce marché diversifié allant des offres grand public aux produits haut de gamme. Notre produit devra se démarquer tout en restant économiquement accessible, notamment pour des partenariats institutionnels, sans compromettre un niveau technologique satisfaisant. 

\subsection{Solution choisie [partie modifiée]}
Notre projet, bien que pouvant s'adresser à un public très large, vise tout particulièrement les personnes en situation d'isolement qui sont attachées à jouer sur un échiquier physique. Nous sommes conscients que notre échiquier interactif intéressera moins les joueurs qui sont habitués à jouer en ligne puisque des applications comme chess.com leur permettent déjà de trouver des adversaires de leur niveau gratuitement. Notre projet répondra donc au besoin suivant: permettre aux joueurs qui sont habitués aux échiquiers physiques d'avoir accès aux mêmes fonctionnalités que celles auxquelles les joueurs en ligne ont accès, en particulier la simulation d'adversité et les indications. Nous ciblons particulièrement les personnes âgées qui ont joué sur des échiquiers toute leur vie et qui n'ont plus la possibilité d'aller trouver des adversaires au club d'échecs du quartier. Dans cette optique, l'échiquier interactif sera simple d'utilisation et les cases indiquées seront illuminées clairement. \\
De plus nous proposerons une version en kit pour aller dans le sens de la tendance DIY. Cela aura pour objectif de viser une autre clientèle, à savoir des joueurs plus jeunes qui aime faire des choses de leurs propres mains. Cette approche ne ferait pas augmenter le coût du produit de base et sera un avantage face à la concurrence. Ainsi nous aurons deux versions du produit, visant deux cibles bien distinctes et donc un marché plus large. 
\textit{Source : \url{https://www.accio.com/business/best-selling-smart-chess-board?utm_source=chatgpt.com}}.
Nous proposerons également dans un second temps plusieurs type de gamme de produits avec des extensions (personnalisation des couleurs, pièces artisanales, matériaux premium) et des versions thématiques (échiquiers inspirés de licences populaires, éditions limitées).  
Ensuite, nous pourrions créer une plateforme communautaire en ligne: une application ou un site dédié pourrait offrir des fonctionnalités comme le partage de parties, un marché de revente pour les kits d’occasion… Nous pensons pour le moment que cette option n'est pas indispensable et apporterait peu de valeur ajoutée par rapport au public visé. 
Finalement, nous espérons intégrer des partenariats éducatifs avec des écoles, des clubs d’échecs voire des maisons de retraite afin de positionner Smart Chess comme un outil pédagogique de référence.

\section{Compétences du groupe}

Notre équipe d’étudiants dispose de compétences variées et complémentaires permettant d’assurer le développement du projet \textit{Smart Chess} :

\subsection{Compétences techniques}
\begin{itemize}
    \item Programmation (Python, C++) et algorithmique appliquée aux jeux d’échecs.
    \item Électronique et systèmes embarqués (capteurs, actionneurs, cartes Raspberry Pi).
    \item Conception mécanique (système de déplacement des pièces, modélisation 3D).
    \item Développement et réglage d’une intelligence artificielle d’échecs.
\end{itemize}

\subsection{Compétences organisationnelles}
\begin{itemize}
    \item Gestion de projet (planification, organigramme, suivi des tâches).
    \item Utilisation d’outils collaboratifs (GitHub, Google Drive, Discord, WhatsApp).
    \item Communication interne et externe (réunions régulières, présentations client).
\end{itemize}

\subsection{Compétences transversales et pédagogiques}
\begin{itemize}
    \item Conception d’outils éducatifs (IA open source, échiquier en kit).
    \item Vulgarisation technique et accompagnement pédagogique.

\end{itemize}

\section{Conclusion}
L'organisation que nous avons définie et les différents outils de travail collaboratif que nous avons mis en place nous permettent d'avancer efficacement dans notre projet. Pour l'instant, les délais prévus dans notre diagramme de Gantt sont respectés, comme indiqué sur notre kanban (en annexe). 




\section{Annexe}

\begin{landscape}
\includegraphics[width=20cm]{kanban.png}   
\end{landscape}

\includepdf[pages=-,landscape=true]{PresentationClient/SlidePresentation.pdf}



\end{document}
